% Generated draft: requires mathematical and source review.
\documentclass[11pt]{article}
\usepackage[margin=25mm]{geometry}
\usepackage{fontspec}
\setmainfont{Latin Modern Roman}
\usepackage{amsmath,amssymb,mathtools,unicode-math}
\setmathfont{Latin Modern Math}
\usepackage{enumitem,needspace,xurl,hyperref}
\setlength{\parindent}{0pt}
\setlength{\parskip}{5pt}
\setlength{\emergencystretch}{3em}
\newcommand{\N}{\mathbb{N}}
\newcommand{\Z}{\mathbb{Z}}
\newcommand{\Q}{\mathbb{Q}}
\newcommand{\R}{\mathbb{R}}
\newcommand{\C}{\mathbb{C}}
\newcommand{\F}{\mathbb{F}}
\newcommand{\A}{\mathbb{A}}
\newcommand{\PP}{\mathbb{P}}
\newcommand{\E}{\mathbb{E}}
\newcommand{\eps}{\varepsilon}
\newcommand{\dd}{\,\mathrm{d}}
\begin{document}
\section*{UnsolvedMath ID: 30001470}
\textbf{Source problem code:} catalogue entry
\subsection{Definitions and mathematical statement}
The supplied record concerns a real matrix \(L\in M_{\ell,m}(\R)\), where \(M_{\ell,m}(\R)\) denotes the space of real \(\ell\)-by-\(m\) matrices. For a vector \(x=(x_1,\ldots,x_r)\) with a positive number of coordinates, write \(\|x\|_\infty=\max_{1\leq i\leq r}|x_i|\). The record does not specify the admissible dimensions \(\ell,m\) or the domain of the parameter \(k\).

The supplied approximation function is
\[
\psi_L(k)=\inf_{\substack{q\in\Z^m,\ p\in\Z^\ell\\ \|q\|_\infty\leq k}}\|Lq-p\|_\infty.
\]
There is an unresolved defect in this formulation: whenever \(k\geq0\), the displayed admissible set includes \(q=0\) and \(p=0\), so the formula as written gives \(\psi_L(k)=0\) for every \(L\). A restriction excluding \(q=0\) could change this, but no such restriction is present in the supplied record, and it is not inserted here.

The source asks which exceptional behaviors of \(\psi_L\), relative to the generic case, are possible and how large the corresponding sets of matrices are. It does not define the generic case, the properties constituting exceptional behavior, or the notion of size of an exceptional set. Thus this is a broad classification question whose intended nontrivial formulation remains unresolved, rather than a fully specified quantified conjecture.
\subsection{Short English statement}
The question asks how matrix Diophantine approximation can differ from typical behavior and how large the exceptional sets are. The supplied formula permits the zero approximation and therefore vanishes for nonnegative \(k\). The intended admissibility restriction, meaning of typical behavior, and notion of exceptional-set size require clarification.
\subsection{Sources}
\begin{itemize}
\item Local UnsolvedMath registry: permanent ID \(30001470\), title ``Exceptional Diophantine Approximation by Matrix Systems,'' source problem code \texttt{catalogue entry}, in the Oberwolfach Reports Open Problems collection. The supplied proposed year is 2010; no individual proposer is identified.
\item Supplied source citation: \emph{Non-positive Curvature and Geometric Structures in Group Theory} (2010), Oberwolfach Reports, DOI 10.4171/owr/2010/20, pp.\ 1165--1224. Source URL and supplied literature reference labeled ``original OWR report'': \url{https://doi.org/10.4171/owr/2010/20}. The background describes these reports as records of mathematical workshops and their open-problem discussions.
\item Supplied literature assessment, checked on \texttt{2026-08-21}: status labeled ``open.'' The assessment describes broad questions about exceptional behavior and exceptional-set size, without a specified theorem or sufficiently narrow class for a definitive resolution, and reports that no general answer was found. This is supplied dated curation derived from the OWR material, not an independently verified literature assessment. Time-sensitive status and cited claims require independent verification.
\item The registry labels the statement \texttt{corrected\_verified} and reports a manual semantic, typographic, and LaTeX review of a meaning-preserving cleanup. These are supplied review claims; they do not resolve the formulation issues identified above. The cited report has not been independently consulted for this draft.
\end{itemize}
\end{document}
